\documentclass{article}
\usepackage[utf8]{inputenc}
\usepackage{graphicx}
\usepackage{amsmath}
\usepackage{amssymb}
\usepackage{url}
\usepackage{tabularx}
\usepackage[colorlinks=true,allcolors=blue]{hyperref}
\usepackage{caption}
\usepackage{pdfpages}
\usepackage{multicol}
\usepackage[ngerman]{babel}
\usepackage{multicol}
\usepackage{geometry}	% [textwidth=17cm,textheight = 24.5cm]
% Um die Dokumentränder zu sehen: [showframe] vor {geometry} einfügen
\usepackage{bbold}
\usepackage{makecell}
\usepackage{float}
\usepackage{braket}
\usepackage{cancel}
\usepackage{braket}
\usepackage{mathtools}
\usepackage{xltabular}
\usepackage{float}
\usepackage{MnSymbol}
\usepackage{multirow}
\usepackage{chemformula}

\renewcommand{\cellalign}{tl}

\setlength{\columnsep}{1cm}
\setlength\parindent{0pt}

\newcommand{\diff}{\mathrm{d}}
\newcommand{\konst}{\mathrm{konst.}}
\newcommand{\lr}[1]{\left(#1\right)}
\newcommand{\lreck}[1]{\left[#1\right]}
\newcommand{\dkon}[3]{\left(\frac{\partial #1}{\partial #2}\right)_{#3}}

\definecolor{gruen}{HTML}{038133}

\newcommand{\commenttext}[1]{\textcolor{gruen}{\text{#1}}}
\newcommand{\comment}[1]{{\color{gruen} {#1}}}

\begin{document}

\begin{center}
	\Large{\textbf{Analysis II: Mehrdimensionale Analysis}}\\
	\small{613348}\\
	\small{\today}
\end{center}

\tableofcontents


\newpage


\section{Differenzierbarkeit}
\subsection{Partielle Differenzierbarkeit}
Die partiellen Ableitungen von $f$ am Punkt $\vec{a}$ sind die Richtungsableitungen entlang der Richtungsvektoren $\vec{e}_i$.
	\begin{equation}
		D_{\vec{e}_i}f(\vec{a}) \equiv \frac{\partial f}{\partial \vec{e}_i}(\vec{a}) \equiv \lim\limits_{h\to 0} \frac{f(\vec{a} + h \vec{e}_i) - f(\vec{a})}{h}
	\end{equation}

{\color{gruen} Merke: Die Richtungsableitung entlang $\vec{v}$ kann man sich aus der 1d Ableitung, die mit $h\to 0$ formuliert ist, ableiten
	\begin{align*}
		f'(a) &= \lim\limits_{h\to 0} \frac{f(a+h) - f(a)}{h} \\
		\Rightarrow D_{\vec{v}}f(\vec{a}) &= \lim\limits_{h\to 0} \frac{f(\vec{a} + h\vec{v}) - f(\vec{a})}{h}
	\end{align*}
}


\subsection{Totale Differenzierbarkeit}
Eine Funktion $\vec{f}$ ist total differenzierbar, wenn sie in alle Richtungen differenzierbar ist. Dann wird der Ableitung am Punkt $\vec{a}$ die \textbf{Jacobi-Matrix} $J_f(a)$
	\begin{equation}
		\diff\vec{f}(\vec{a}) \equiv J_{\vec{f}}(\vec{a}) =
		\begin{pmatrix}
			\frac{\partial f_1}{\partial x_1}(a) & \dots & \frac{\partial f_1}{\partial x_n}(a) \\
			\vdots & \ddots & \vdots \\
			\frac{\partial f_m}{\partial x_1}(a) & \dots & \frac{\partial f_m}{\partial x_n}(a)
		\end{pmatrix}
	\end{equation}
zugeordnet.\\
\\
Die Änderung der Funktion $\vec{f}$, wenn man in die Richtung $\vec{v}$ schreitet ergibt sich dann aus
	\begin{equation}
		J_{\vec{f}}(\vec{a})\vec{v} =
		\begin{pmatrix}
			\frac{\partial f_1}{\partial x_1}(\vec{a}) & \dots & \frac{\partial f_1}{\partial x_n}(\vec{a}) \\
			\vdots & \ddots & \vdots \\
			\frac{\partial f_m}{\partial x_1}(\vec{a}) & \dots & \frac{\partial f_m}{\partial x_n}(\vec{a})
		\end{pmatrix}
		\cdot
		\begin{pmatrix}
			v_1\\
			\vdots\\
			v_n
		\end{pmatrix}
	\end{equation}

Im Falle einer skalaren Funktion $f:\mathbb{R}^3\to\mathbb{R}$ und einem infinitesimal kleinen Schritt $\diff \vec{x}$ ergibt sich das sogenannte \textbf{totale Differential}
	\begin{equation}
		\diff f := J_f(\vec{a})\cdot\diff\vec{x} = \frac{\partial f}{\partial x_1} \diff x_1 + \frac{\partial f}{\partial x_2} \diff x_2 + \frac{\partial f}{\partial x_3} \diff x_3 = \sum_{i=1}^3 \frac{\partial f}{\partial x_i} \diff x_i
	\end{equation}

Im Falle einer vektoriellen Funktion $f:\mathbb{R}^3\to\mathbb{R}^3$ und einem infinitesimal kleinen Schritt $\diff \vec{x}$ ergibt sich das \textbf{totale Differential}
	\begin{equation}
		\diff\vec{f} := J_{\vec{f}}(\vec{a})\cdot\diff\vec{x} =
		\begin{pmatrix}
			\frac{\partial f_1}{\partial x_1}(\vec{a}) & \frac{\partial f_1}{\partial x_2}(\vec{a}) & \frac{\partial f_1}{\partial x_3}(\vec{a}) \\
			\frac{\partial f_2}{\partial x_1}(\vec{a}) & \frac{\partial f_2}{\partial x_2}(\vec{a}) & \frac{\partial f_2}{\partial x_3}(\vec{a}) \\
			\frac{\partial f_3}{\partial x_1}(\vec{a}) & \frac{\partial f_3}{\partial x_2}(\vec{a}) & \frac{\partial f_3}{\partial x_3}(\vec{a})
		\end{pmatrix}
		\cdot
		\begin{pmatrix}
			\diff x_1 \\ \diff x_2 \\ \diff x_3
		\end{pmatrix}
		=
		\begin{pmatrix}
			\sum_{i=1}^3 \frac{\partial f_1}{\partial x_i} \diff x_i\\
			\sum_{i=1}^3 \frac{\partial f_2}{\partial x_i} \diff x_i\\
			\sum_{i=1}^3 \frac{\partial f_3}{\partial x_i} \diff x_i
		\end{pmatrix}
	\end{equation}



%Eine Funktion $f:\mathbb{R}^n\to\mathbb{R}^m$ ist in $a\in\Omega\subset \mathbb{R}^n$ differenzierbar, wenn sie durch $f(x) \approx f(a) + L(x+a)$ approximiert werden kann, d.h. für den Restterm $R(x-a) = f(x) - f(a) - L(x-a)$ gilt
%	\begin{equation}
%		\lim\limits_{x\to a} \frac{||R(x-a)||}{||x-a||} \stackrel{!}{=} 0
%	\end{equation}
%
%Die \textbf{Ableitung von $f$ am Punkt $a$} ist die \underline{Abbildung} $L \equiv \diff f(a) \equiv f'(a)$.\\


\newpage

{\color{blue}
\hrule
	\fbox{Beispiel}\\
	geg: $f:\mathbb{R}^2\to \mathbb{R}^3,\ (x,y)\mapsto (x^2+y^2, x, y)^T$\\
	\\
	Die Ableitungen für die Jacobi-Matrix (die Ableitung von $f$) sind
		$$\frac{\partial f_1}{\partial x} = 2x,\
			\frac{\partial f_2}{\partial x} = 1,\
			\frac{\partial f_3}{\partial x} = 0,\
			\frac{\partial f_1}{\partial y} = 2y,\
			\frac{\partial f_2}{\partial y} = 0,\
			\frac{\partial f_3}{\partial y} = 1$$
	
	Somit ist die Jacobi-Matrix
		$$J_f(a) =
		\begin{pmatrix}
			\frac{\partial f_1}{\partial x} & \frac{\partial f_1}{\partial y} \\[0.5em]
			\frac{\partial f_2}{\partial x} & \frac{\partial f_2}{\partial y} \\[0.5em]
			\frac{\partial f_3}{\partial x} & \frac{\partial f_3}{\partial y}
		\end{pmatrix}
		=
		\begin{pmatrix}
			2x & 2y \\
			1 & 0 \\
			0 & 1
		\end{pmatrix}
		$$
\hrule
}


\subsection{Rechenregeln}
	\begin{itemize}
		\item \textbf{Kettenregel}
			\begin{equation}
				\diff(f\circ g)(x) = \diff f(g(x)) \circ \diff g(x)
			\end{equation}

		\item \textbf{Produktregel}
			\begin{equation}
				\diff(f\cdot g)(x) = \diff f(x)\cdot g(x) + f(x)\cdot \diff g(x)
			\end{equation}
		
		\item \textbf{Quotientenregel}
			\begin{equation}
				\diff\lr{ \frac{f}{g} }(x) = \frac{f'(x)\cdot g(x) - f(x)\cdot g'(x)}{g(x)^2}
			\end{equation}
	\end{itemize}

\vspace{0.5cm}

{\color{blue}
\hrule
	\fbox{Beispiel}\\
	geg: $g:[0,T]\to \mathbb{R}^3$, $f:\mathbb{R}^3\to \mathbb{R}$, $g(t) = (\cos t, \sin t, t)^T$, $f(x,y,z) = x + yz$\\
	\\
	Berechne die Ableitungen, d.h. die Jacobi-Matrizen für $f,g$.
		\begin{align*}
			\diff f(x,y,z) &=
			\begin{pmatrix}
				\frac{\partial f}{\partial x}(x,y,z) \\
				\frac{\partial f}{\partial y}(x,y,z) \\
				\frac{\partial f}{\partial z}(x,y,z)
			\end{pmatrix}
			=
			\begin{pmatrix}
				1 \\
				z \\
				y
			\end{pmatrix} \\
			\diff g(t) &=
			\begin{pmatrix}
				\frac{\partial g_1}{\partial t}(t) \\
				\frac{\partial g_2}{\partial t}(t) \\
				\frac{\partial g_3}{\partial t}(t)
			\end{pmatrix}
			=
			\begin{pmatrix}
				-\sin t\\
				\cos t\\
				1
			\end{pmatrix} \\
			\Rightarrow \diff(f\circ g)(t) &= \diff f(g(t)) \circ g(t) =
			\begin{pmatrix}
				1\\
				t\\
				\sin t
			\end{pmatrix}
			\cdot
			\begin{pmatrix}
				-\sin t\\
				\cos t\\
				1
			\end{pmatrix}
			=
			\cancel{-\sin t} + t\cos t + \cancel{\sin t} = t\cos t
		\end{align*}
\hrule
}


\newpage


\subsection{Taylor-Entwicklung}
Das Taylor-Polynom einer Funktion $f(x)$ am Entwicklungpunkt $h$ \textcolor{gruen}{(wobei $x,h$ Vektoren sind)} bis zur $k$-ten Ordnung, ist
	\begin{align}
		T_{f,h,k}(x) &= f(h) + \sum_{j=1}^k \frac{1}{j!} D^{(j)}f(h)(x-h)^j = \\
		&= f(h) + D^{(1)}f(h)(x-h) + \frac{1}{2} \cdot D^{(2)}f(h)(x-h,x-h) + ... = \\
		&= f(h) + D^{(1)}f(h)\cdot (x-h) + \frac{1}{2} \cdot  (x-h)^T \cdot D^{(2)}f(h) \cdot (x-h) + ...
	\end{align}

\vspace{0.5cm}

{\color{blue}
\hrule
	\fbox{Beispiel}\\
	ges: $T_{f,(0,0),2}(x,y)$ mit $f(x,y) = \exp(x\cdot y)$
		\begin{align*}
			T_{f,(0,0),2}(x,y) &= f(0,0) + D^{(1)}f(0,0)\cdot
			\begin{pmatrix}
				x\\y
			\end{pmatrix}
			 + 
			 \frac{1}{2} \cdot
			 \begin{pmatrix}
				x & y
			\end{pmatrix}
			\cdot
			D^{(2)}f(0,0)\cdot
			\begin{pmatrix}
				x\\y
			\end{pmatrix}
		\end{align*}
	
	Berechne die Ableitungen mittels Jacobi-Determinante.
		\begin{align*}
			D^{(1)}f(0,0) &=
			\begin{pmatrix}
				\partial_x f(0,0) & \partial_y f(0,0)
			\end{pmatrix}
			=
			\left.
			\begin{pmatrix}
				y\exp(xy) & x\exp(xy)
			\end{pmatrix}
			\right|_{(0,0)}
			=
			\begin{pmatrix}
				0 & 0
			\end{pmatrix} \\
			\Rightarrow D^{(2)}f(0,0) &=
			D^{(1)}
			\begin{pmatrix}
				\partial_x f(0,0) & \partial_y f(0,0)
			\end{pmatrix}
			=
			\begin{pmatrix}
				\partial_x^2f(0,0) & \partial_y \partial_x f(0,0) \\
				\partial_y \partial_x f(0,0) & \partial_y^2f(0,0)
			\end{pmatrix}
			= \\
			&=
			\left.
			\begin{pmatrix}
				\partial_x y\exp(xy) & \partial_y y\exp(xy) \\
				\partial_y x\exp(xy) & \partial_y x\exp(xy)		
			\end{pmatrix}
			\right|_{(0,0)} =
			\left.
			\begin{pmatrix}
				y^2\exp(xy) & (1+xy)\exp(xy) \\
				(1+xy)\exp(xy) & x^2\exp(xy)		
			\end{pmatrix}
			\right|_{(0,0)} =\\
			&=
			\begin{pmatrix}
				0 & 1 \\
				1 & 0
			\end{pmatrix}
		\end{align*}

	Somit wird das Taylor-Polynom
		\begin{align*}
			T_{f,(0,0),2}(x,y) &= 1 +
			\begin{pmatrix}
				0 & 0
			\end{pmatrix}
			\begin{pmatrix}
				x\\ y
			\end{pmatrix}
			+
			\frac{1}{2}
			\begin{pmatrix}
				x & y
			\end{pmatrix}
			\begin{pmatrix}
				0 & 1\\
				1 & 0
			\end{pmatrix}
			\begin{pmatrix}
				x\\y
			\end{pmatrix}
			= 1 + 0 + \frac{1}{2}
			\begin{pmatrix}
				x & y
			\end{pmatrix}
			\begin{pmatrix}
				y\\x
			\end{pmatrix}
			=
			1 + xy
		\end{align*}

	\begin{figure}[H]
		\centering
		\includegraphics[width=0.45\textwidth]{Grafiken/Ana II Mehrdimensionale Taylor-Entwicklung.pdf}
	\end{figure}
\hrule
}


\newpage


\subsection{Lokale Extrema von Funktionen $f:\Omega\subset\mathbb{R}^n\to \mathbb{R}$}
Ein Punkt $a$ ist ein lokales Minimum (Maximum) [Sattelpunkt], wenn
	\begin{enumerate}
		\item Der Gradient verschwindet
			\begin{equation}
				f'(a) = 0
			\end{equation}
		
		\item Die Hesse-Matrix nur positive (negative) [positive und negative] Eigenwerte hat oder äquivalent
			\begin{equation}
				\forall x\in\Omega:\quad
				\braket{x,Hf(...)x} 
				\begin{cases}
				> 0, & \text{Minimum} \\
				< 0, & \text{Maximum} \\
				\text{sonst}, & \text{Sattelpunkt}
				\end{cases}
			\end{equation}
	\end{enumerate}

\vspace{0.5cm}

{\color{blue}
\hrule
	\fbox{Beispiel}\\
	geg: $f(x,y) = x^2+y^2+xy+1$\\
	\\
	Finde zunächst die kritischen Punkte. Berechne also den Gradienten und setze ihn null.
		\begin{align*}
			f'(x,y) &=
			\begin{pmatrix}
				\partial_x f(x,y)\\
				\partial_y f(x,y)
			\end{pmatrix}
			=
			\begin{pmatrix}
				2x + y\\
				2y + x
			\end{pmatrix}
			\stackrel{!}{=}
			\begin{pmatrix}
				0\\
				0
			\end{pmatrix}
		\end{align*}
	
	Dies wird erfüllt für den Punkt $a = (0,0)$.\\
	Berechne also die Hesse-Matrix für diesen Punkt.
		\begin{align*}
			Hf(x,y) &=
			\begin{pmatrix}
				\partial_x^2f(0,0) & \partial_x\partial_y f(x,y)\\
				\partial_y^2f(0,0) & \partial_y\partial_x f(x,y)		
			\end{pmatrix}
			=
			\begin{pmatrix}
				2 & 1\\
				1 & 2
			\end{pmatrix}\\
			\Rightarrow \braket{x, Hf(x,y)x} &=
			\begin{pmatrix}
				x & y
			\end{pmatrix}
			\begin{pmatrix}
				2 & 1\\
				1 & 2
			\end{pmatrix}
			\begin{pmatrix}
				x\\
				y
			\end{pmatrix}
			=
			\begin{pmatrix}
				x & y
			\end{pmatrix}
			\begin{pmatrix}
				2x + y\\
				2y + x
			\end{pmatrix}
			=
			2x^2 + yx + xy + 2y^2 = \\
			&= (x + y)^2 + x^2 + y^2 > 0
		\end{align*}
	
	Somit ist $a=(0,0)$ ein lokales Minimum.
	
	\begin{figure}[H]
		\centering
		\includegraphics[width=0.45\textwidth]{Grafiken/Ana II Lokale Extrema.pdf}
	\end{figure}
\hrule
}


\newpage


\subsection{Lokale Extrema unter Nebenbedingungen}
Betrachte eine Funktion $g$, die unter $m$ Nebenbedingungen $f_i$ extremal, d.h. minimal oder maximal, werden soll.\\
\\
Die hinreichende Bedingung für Extremalität ist
	\begin{enumerate}
		\item Verschwinden des Gradienten (notwendige Bedingung)
			\begin{equation}
				\nabla G := \nabla g - \sum_{i=1}^m \lambda_i \nabla f_i = 0
			\end{equation}
		
		\item Für alle Vektoren $v$, die senkrecht auf dem Gradienten an der Extremalstelle stehen, gilt
			\begin{equation}
				\braket{v,HG(z^\star) v}
				\begin{cases}
					> 0, & \text{Minimum}\\
					< 0, & \text{Maximum}
				\end{cases}
			\end{equation}
	\end{enumerate}

\textbf{Anleitung}
	\begin{enumerate}
		\item \textbf{Funktionen aufstellen}\\
		Stelle die zu minimierende (maximierende) Funktion $g$, die Nebenbedingungen $f_i$, sowie die Funktion
			\begin{equation}
				G = g - \sum_{i=1}^m \lambda_i f_i
			\end{equation}
		auf.


		\item \textbf{Notwendige Bedingung für Extremalität}\\
		Setze $g$ und die $f_i$ in die Gl.
			\begin{equation}
				\nabla G = \nabla g - \sum_{i=1}^m \lambda_i \nabla f_i = 0
			\end{equation}

		 ein und gewinne daraus durch Elimination der Lagrange-Multiplikatoren $\lambda_i$ und Verwendung der Nebenbedingung(en)
		 	\begin{equation}
				 f_i = 0
		 	\end{equation}

		die lokale(n) Extremstelle(n) $z^\star = (x^\star,y^\star,...)$.


		\item \textbf{Lagrange-Multiplikatoren}\\
		Bestimme die zu den Extremalstellen $z^\star$ gehörenden Lagrange-Multiplikatoren.


		\item \textbf{Senkrechte Vektoren zu $\nabla f_i$}\\
		Bestimme die zum Gradienten der Nebenbedingungen $f_i$ an der Extremalstelle $z^\star$ senkrechten Vektoren $v$ über die Bedingung
			\begin{align*}
				\braket{v,\nabla f(z^\star)} = 0
			\end{align*}	


		\item \textbf{Hesse-Matrix}\\
		Bestimme die Hesse-Matrix
			\begin{equation}
				HG(z^\star)
			\end{equation}
		
		und berechne
			\begin{equation}
				\braket{v,HG(z^\star) v}
			\end{equation}
	\end{enumerate}


\newpage


{\color{blue}
\hrule
	\fbox{Beispiel}\\
	geg: $g(x,y) = xy + 14$, $f(x,y) = x^2 + y^2 - 18$\\
	\\
	Stelle zunächst $G(x,y)$ auf
		$$G(x,y) = xy + 14 - \lambda(x^2 + y^2 - 18)$$
	
	Bestimme nun die Extremalstellen über 
		\begin{align*}
			\nabla G(x,y) &=
			\begin{pmatrix}
				\partial_x G(x,y)\\
				\partial_y G(x,y)
			\end{pmatrix}
			=
			\begin{pmatrix}
				y - 2\lambda x\\
				x - 2\lambda y
			\end{pmatrix}
			\stackrel{!}{=}
			\begin{pmatrix}
				0\\
				0
			\end{pmatrix}
		\end{align*}
	
	Eliminiere $\lambda$
		\begin{align*}
			\begin{cases}
				y - 2\lambda x &= 0\\
				x - 2\lambda y &= 0
			\end{cases}
			\Leftrightarrow
			\begin{cases}
				\lambda &= \frac{y}{2x}\\
				\lambda &= \frac{x}{2y}
			\end{cases}
			\Leftrightarrow
			\frac{y}{2x} = \frac{x}{2y}
			\Leftrightarrow
			x^2 = y^2
		\end{align*}
	
	Verwende nun die Nebenbedingung $f=0$, um die Extremalstellen final zu bestimmen.
		\begin{align*}
			x^2 + y^2 - 18 &= 0 \Leftrightarrow x^2 + x^2 -18 = 0 \Leftrightarrow x = \pm 3
		\end{align*}
	
	Somit sind die Extremalstellen $z_1^\star = (3,3), z_2^\star = (-3,3), z_3^\star = (3,-3), z_4^\star = (-3,-3)$.\\
	Die dazugehörigen Lagrange-Multiplikatoren sind
		\begin{align*}
			\lambda_1 = \lambda_4 &= \frac{3}{2\cdot 3} = \frac{1}{2}\\
			\lambda_2 = \lambda_3 &= -\frac{3}{2\cdot 3} = -\frac{1}{2}
		\end{align*}

	Bestimme nun die zu den Gradienten an den Extremalstellen senkrechten Vektoren.
		\begin{align*}
			\bullet \braket{v, \nabla f(3,3)} =
			\begin{pmatrix}
				v_1 & v_2
			\end{pmatrix}
			\cdot
			\left.
			\begin{pmatrix}
				2x\\
				2y
			\end{pmatrix}
			\right|_{(\pm 3,\pm 3)}
			=
			\pm 6v_1 + \pm 6v_2 \stackrel{!}{=} 0 \\
			\Rightarrow v_1 &= v_2\\
			\bullet \braket{v, \nabla f(3,3)} =
			\begin{pmatrix}
				v_1 & -v_2
			\end{pmatrix}
			\cdot
			\left.
			\begin{pmatrix}
				2x\\
				2y
			\end{pmatrix}
			\right|_{(\pm 3,\mp 3)}
			=
			\pm 6v_1 \mp 6v_2 \stackrel{!}{=} 0\\
			\Rightarrow v_1 &= v_2
		\end{align*}
	
	
	Bestimme die Hesse-Matrix
		\begin{align*}
			HG(x,y) &=
			\begin{pmatrix}
				\partial_x^2 G(x,y) & \partial_x\partial_y G(x,y) \\
				\partial_y\partial_x G(x,y)  & \partial_y^2 G(x,y) 
			\end{pmatrix}
			=
			\begin{pmatrix}
				-2\lambda & 1 \\
				1 & -2\lambda
			\end{pmatrix}
		\end{align*}
	
	Bestimme, ob die Extremalstellen Minima oder Maxima sind.
		\begin{align*}
			\bullet \braket{v, HG(3,3) v} = \braket{v, HG(-3,-3) v} &=
			\begin{pmatrix}
				v_1 & -v_1
			\end{pmatrix}
			\cdot
			\begin{pmatrix}
				-2\frac{1}{2} & 1\\
				1 & -2\frac{1}{2}
			\end{pmatrix}
			=
			-2v_1^2 < 0\\
			\bullet \braket{v, HG(-3,3) v} = \braket{v, HG(3,-3) v} &=
			\begin{pmatrix}
				v_1 & v_1
			\end{pmatrix}
			\cdot
			\begin{pmatrix}
				-2\lr{-\frac{1}{2}} & 1\\
				1 & -2\lr{-\frac{1}{2}}
			\end{pmatrix}
			=
			2v_1^2 > 0
		\end{align*}
	
	Somit sind $(3,3),(-3,-3)$ Maxima, während $(-3,3),(3,-3)$ Minima.
	

	\begin{figure}[H]
		\centering
		\includegraphics[width=0.75\textwidth]{Grafiken/Ana II Lokale Extrema mit Nebenbedingungen.pdf}
	\end{figure}
\hrule
}



\newpage


\section{Integralrechnung in $\mathbb{R}^n$}
\subsection{Rechenregeln}
	\begin{itemize}
		\item \textbf{Vertauschung von Integration und Differentiation}\\
		Für überall stetige und in $y$ partiell differenzierbare Funktionen $f$ gilt
			\begin{equation}
				\frac{\partial}{\partial y_k} \int_Q f(x,y)\, \diff x\diff y = \int_Q \frac{\partial f}{\partial y_k}(x,y)\, \diff x\diff y
			\end{equation}
		
		
		\item \textbf{Satz von Fubini (Vertauschung von Integralen)}\\
		Die Integrationsreihenfolge darf bei der Integration über kompakte Quader \textcolor{gruen}{(d.h. abgeschlossen und beschränkte Quader, z.B. $[a,b], [c,d]$)} vertauscht werden darf
			\begin{align}
				\int_{[a,b]\times [c,d]} f(x,y)\, \diff x\diff y &= \int_a^b \lr{ \int_c^d f(x,y)\,\diff y }\diff x = \int_c^d \lr{ \int_a^b f(x,y)\,\diff x }\diff y
			\intertext{bzw. allgemeiner für zwei Quader $Q_1\in\mathbb{R}^p, Q_2\in\mathbb{R}^q$}
				\int_{Q_1\times Q_2} f(x,y)\, \diff^{p\times q} (x,y) &= \int_{Q_1} \lr{ \int_{Q_2} f(x,y)\,\diff^q y }\diff^p x = \int_{Q_2} \lr{ \int_{Q_1} f(x,y)\,\diff^p  }\diff^q y 
			\end{align}
		
		
		\item \textbf{Transformationssatz}\\
		Sei $\varphi: U\to V, u\mapsto \varphi(u)$ eine Koordinatentransformation und $f:V$ eine stetige Funktion
			\begin{equation}
				\int_V f(x) \diff^n x = \int_U f(\varphi(u)) \cdot |\det J_\varphi(u)|\, \diff^n u
			\end{equation}
	\end{itemize}

	\textcolor{gruen}{Notation: $U$: Vektorraum der substituierten Variablen mit Elementen $u$; $V$: Vektorraum der Koordinaten zu Beginn mit Elementen $x$; $\varphi$: Funktion, die die substituierten Elemente auf die Elemente zu Beginn abbildet.}

\vspace{0.5cm}

{\color{blue}
\hrule
	\fbox{Beispiel}\\
	ges: $\int\ln(x^2+y^2)\,\diff x\diff y$ mit $a\leq \sqrt{x^2 + y^2} \leq b$ und $x,y\geq 0$\\
	\\
	Das Integral entspricht der Integration über eine Kreisscheibe.\\
	Gehe daher in Zylinderkoordinaten über.\\
	Die Koordinatentrafo ist $\varphi:U = [a,b]\times[0,\frac{\pi}{2}]\to \mathbb{R}^2,
	u
	=
	\begin{pmatrix}
		r\\ \theta
	\end{pmatrix}	
	\mapsto
	\varphi\lr{
	\begin{pmatrix}
		r\\ \theta
	\end{pmatrix}}
	=
	\begin{pmatrix}
		x\\ y
	\end{pmatrix}
	=
	\begin{pmatrix}
		r\cos\theta \\
		r\sin\theta
	\end{pmatrix}$.\\
	\textcolor{gruen}{$U$ ergibt sich aus den Beschränkungen für $x,y$.}
		\begin{align*}
			\int\ln(x^2+y^2)\,\diff x\diff y &= \int_a^b \lr{ \int_0^{\pi/2} \ln(r^2)\cdot
			\begin{vmatrix}
				\frac{\partial\varphi_1}{\partial r} & \frac{\partial\varphi_1}{\partial \theta} \\
				\frac{\partial\varphi_2}{\partial r} & \frac{\partial\varphi_2}{\partial \theta}
			\end{vmatrix}
			\,\diff \theta
			}
			\diff r
			= \frac{\pi}{2} \int_a^b 2\ln r \cdot
			\begin{vmatrix}
				\cos\theta & -r\sin\theta \\
				\cos\theta & r\cos\theta
			\end{vmatrix}
			\,\diff r = \\
			&= \pi \int_a^b \ln r\cdot r\,\diff r = ... = \pi \lr{ \frac{b^2\ln b - a^2\ln a}{2} - \frac{b^2 - a^2}{4}}
		\end{align*}


\newpage


	\begin{figure}[H]
		\centering
		\includegraphics[width=0.5\textwidth]{Grafiken/Ana II Transformationssatz.pdf}
	\end{figure}

\hrule
}


\subsection{Volumen zwischen zwei Funktionen}
Hängen die Integralgrenzen einer Variablen von einer anderen ab, so wird zunächst über die erste Variable integriert.\\
Integration über die Menge $B = \{(x,y)\in D\times \mathbb{R}, h_-(x) \leq y \leq h_+(x)\}$ erfolgt über
	\begin{equation}
		\int_B f(x,y)\,\underbrace{\diff x}_{=\diff x_1\cdot ...\cdot \diff x_{n-1}}\diff y = \int_D \lr{ \int_{h_-(x)}^{h_+(x)} f(x,y),\diff y } \diff x
	\end{equation}

\vspace{0.5cm}

{\color{blue}
\hrule
	\fbox{Beispiel}\\
	



\hrule
}


\newpage


\subsection{Volumen- und Oberflächenberechnung von Rotationskörpern im $\mathbb{R}^3$}
Das Volumen eines Rotationskörpers, der durch die Rotation der Funktion $r:[a,b]\to \mathbb{R}, x\mapsto r(x)$ um die $x$-Achse erzeugt wird, ist
	\begin{equation}
		\text{vol}(K) = \pi \int_a^b r(x)^2\,\diff x
	\end{equation}

Die Oberfläche eines Rotationskörpers ergibt sich aus
	\begin{equation}
		\text{surf}(K) = 2\pi \int_a^b r(x)\cdot \sqrt{1 + r'(x)^2}\,\diff x
	\end{equation}

\vspace{0.5cm}

{\color{blue}
\hrule
	\fbox{Beispiel}\\
	Berechne den Rotationskörper, der durch die Rotation der Fläche $\{(x,y,z)\in \mathbb{R}^3, x\in[0,8], y\in[0,\sqrt[3]{x}], z=0\}$ um die $x$-Achse entsteht.\\
	\\
	Für die zu rotierende Funktion gilt $r(x) = y = \sqrt[3]{x}$ und dessen Grenzen $r:[0,8]\to\mathbb{R}$.
		\begin{align*}
			\text{vol}(K) &= \pi \int_0^8 (\sqrt[3]{x})^2\,\diff x = \pi \int_0^8 x^{3/2}\,\diff x = \pi \frac{3}{5} \left. x^{\frac{5}{3}} \right|_0^8 = \pi \frac{3}{5} 8^{\frac{5}{3}} = \pi \frac{3}{5} (2^3)^{\frac{5}{3}} = \\
			&= \pi \frac{3}{5} 2^5 = \frac{96\pi}{5}
		\end{align*}
		
	\begin{figure}[H]
		\centering
		\includegraphics[width=0.6\textwidth]{Grafiken/Ana II Rotationskörper.pdf}	
	\end{figure}
\hrule
}


\newpage


{\color{violet}
\hrule
	\fbox{Herleitung}\\
	zz: $\text{vol}(K) = \pi \int_a^b r(x)^2\,\diff x$
	\begin{figure}[H]
		\centering
		\includegraphics[width=0.8\textwidth]{Grafiken/Ana II Rotationskörper Herleitung.pdf}
	\end{figure}

	Berechne die Querschnittsfläche $\tilde{A}$ mittels ''Berechnung der Fläche zwischen 2 Funktionen'' \textcolor{gruen}{(siehe vorheriges Unterkapitel)}.\\
	Finde dafür zunächst die beiden Funktionen $z_{1,2}(y) \equiv h_{\pm}(y)$. Diese ergeben sich über den Satz des Pythagoras für den Kreisradius
		\begin{align*}
			r(x)^2 &= y^2 + z^2 \qquad\Leftrightarrow\qquad z_{1,2} = \pm \sqrt{r(x)^2 - y^2} \equiv h_\pm(y)
		\end{align*}
	
	$y$ muss einfach von -$r(x)$ bis $r(x)$ integriert werden. Somit ergibt sich
		\begin{align*}
			\tilde{A} &= \int_{\tilde{A}} 1\,\diff y\diff z = \int_{-r(x)}^{r(x)} \lr{ \int_{-\sqrt{r(x)^2-y^2}}^{+\sqrt{r(x)^2-y^2}} 1\,\diff z}\diff y = 2\int_{-r(x)}^{r(x)} \sqrt{r(x)^2 - y^2}\,\diff y \\
			&= 2\int_{-r(x)}^{r(x)} r \sqrt{1- \lr{ \frac{y}{r} }^2} =
		\intertext{Löse dieses Integral mittels der trigonometrischen Substitution $\sin(u) = \frac{y}{r}$, um $\cos(u) = \sqrt{1 - \sin^2(u)}$ verwenden zu können.}
			&= 2r \int_{\arcsin\lr{-\frac{r}{r}}}^{\arcsin\lr{\frac{r}{r}}} \underbrace{\sqrt{1 - \sin^2(u)}}_{=\cos u}\,(\diff u\cdot \cos u\cdot r) = 2r^2 \int_{-\pi/2}^{\pi/2} \cos^2u\,\diff u = \\
			&= 2r^2 \int_{-\pi/2}^{\pi/2} \frac{1}{2}\lr{ \cos(2u) + 1 }\,\diff u = r^2 \lr{ \cancel{\left.\frac{1}{2}\sin(2u) \right|_{-\pi/2}^{\pi/2}} + \lr{\pi/2 - (-\pi/2} } = \\
			&= \pi r(x)^2
		\intertext{Es verbleibt alle Querschnitte von $a$ bis $b$ zu integrieren, was die gesuchte Formel ergibt.}
			\text{vol}(K) &= \pi \int_a^b r(x)^2\,\diff x \qquad\blacksquare
		\end{align*}

\hrule
}


\newpage


\section{Vektoranalysis}
\subsection{Kurvenintegrale}
\subsubsection{Länge einer Kurve}
Die Länge einer Kurve $C$ lässt sich über die zugehörige Parametrisierung $\alpha:[a,b]\to\mathbb{R}^n,t\mapsto \alpha(t)$ durch die folgende Formel berechnen
	\begin{equation}
		L(C) = \int_a^b ||\dot{\alpha}(t)||\,\diff t
	\end{equation}

\textcolor{gruen}{Anschaulich: $s=v\cdot t$ mit $v = ||\dot{a}(t)||$.}


\subsubsection{Kurvenintegral von skalaren Funktionen}
Das skalare Kurvenintegral ist definiert als
	\begin{equation}
		\int_C f\,\diff s = \int_a^b f(\alpha(t))\cdot ||\dot{a}(t)||\,\diff t
	\end{equation}

mit dem sog. skalaren Bogenelement $\diff s$.

\vspace{0.5cm}

{\color{blue}
\hrule
	\fbox{Beispiel}\\
	geg: $f(x,y) = 1 + \frac{y}{3}$, $\alpha(t) = (-3\cos^2(t), 3\sin(t)^2)$
		\begin{align*}
			\int_C f\,\diff s &= \int_0^{2\pi} \lr{ 1 + \frac{3\sin^3(t)}{3} } \cdot
		||
		\begin{pmatrix}
			-9\cos^2(t)\sin(t)\\
			9\sin^2(t)\cos(t)
		\end{pmatrix}
		||\,\diff t
		= \int_0^{2\pi} \lr{ 1 + \frac{3\sin^3(t)}{3} } 9\sin t\cos t = ... = \frac{63}{5}
		\end{align*}
	
	\begin{figure}[H]
		\centering
		\includegraphics[width=0.75\textwidth]{Grafiken/Ana II skalares Kurvenintegral.pdf}
	\end{figure}
\hrule
}





\subsubsection{Kurvenintegral von vektoriellen Funktionen}
Das vektorwertige Kurvenintegral ist definiert als
	\begin{equation}
		\int_C f\,\diff s := \int_a^b \braket{f(\alpha(t)),\dot{\alpha}(t)}\,\diff t
	\end{equation}

\vspace{0.5cm}

{\color{blue}
\hrule
	\fbox{Beispiel}\\
	geg: Vektorfeld $K(x) = -x/||x||^3$, Parametrisierung $\alpha(t)$

	\begin{figure}[H]
		\centering
		\includegraphics[width=0.75\textwidth]{Grafiken/Ana II vektorielles Kurvenintegral.pdf}
	\end{figure}
\hrule
}

\vspace{0.5cm}

Konservative Vektorfelder erfüllen
	\begin{equation}
		\int_\gamma f\,\diff s = 0
	\end{equation}


\subsection{Potentiale von Vektorfeldern}
\subsubsection{•}







\newpage


\section{Stichwortverzeichnis}
\begin{xltabular}{\linewidth}{p{6cm}|X}
	\hline\hline
		Begriff & Erklärung \\
	\hline
		Parametrisierung & Abbildung, die ein Intervall auf die Punkte einer Kurve abbildet
			$$\alpha: I = [\bullet,\bullet] \to \mathbb{R}^n$$
		{\color{gruen}
		Beispiel: Archimedische Spirale
			$$\alpha:[0,T] \to \mathbb{R}^2, t\mapsto
			\begin{pmatrix}
				t\cos t\\
				t\sin t
			\end{pmatrix}
			$$}
		
		\includegraphics[width=0.5\textwidth]{Grafiken/Ana II Archimedische Spirale.pdf} \\
		\hline
			regulär & Eine Kurve $\alpha:I\to\mathbb{R}^n, t\mapsto\alpha(t)$ ist regulär, wenn
				$$\dot{\alpha}(t) \neq 0$$ \\
	\hline\hline
\end{xltabular}

\end{document}
